% ***************************************************
% Conclusion
% ***************************************************
\chapter[Conclusion]{Conclusion}
\label{Chap:Conclusion}

The aim of this thesis was to investigate the production mechanisms behind the diffuse emission of minihalos. A key observational distinction between hadronic production and turbulent reacceleration lies in the high-frequency spectral steepening predicted by the reacceleration scenario. To investigate this, radio observations at 10 GHz for a sample of twelve galaxy cluster minihalos were analysed during this thesis. The results indicate that high-frequency spectral steepening $<10$ GHz is not common among minihalos for which robust flux measurements were obtained.

These datasets were calibrated and jointly imaged to combine the higher angular resolution of C-configuration observations with the better sensitivity to extended emission of D-configuration data. A custom calibration stage was implemented to improve the removal of radio-frequency interference within the observations. The resulting images were source-subtracted to isolate the diffuse emission of the minihalo, and subsequently utilised to calculate the total flux density. This was successfully completed for four clusters, RXJ1720+2638, A478, RXCJ1115+0129, and RXJ2129+0005. The majority of the remaining clusters displayed evidence of a minihalo, however, imaging artefacts prevented a reliable measurement of the total minihalo emission.

The resulting flux densities were compared to previous studies and used to calculate the spectral indices for each minihalo. Three of these clusters (RXJ1720+2638, RXCJ1115+0129, RXJ2129+0005) are consistent with a single power-law up to 10 GHz, with spectral indices in the range $\alpha\approx1$--1.1. One cluster (A478) exhibits potential steepening, with a derived spectral index of $\alpha=2.02$, compared to the lower frequency upper bound of $\alpha\leq1$. The single power-law spectrum is broadly consistent with the hadronic scenario, while the steepening is more consistent with the turbulent reacceleration method. These results suggest that both mechanisms can operate within the minihalo, with hadronic production potentially playing a dominant role in many systems. 

The measured spectral indices can be biased by the defined spatial mask, \textit{uv}-coverage, and source-subtraction between observations at different frequencies. Differences in masking strategies and \textit{uv}-coverage between studies makes comparison to the results here more ambiguous, as spatial filtering from mismatched coverages can artificially reduce the measured flux. Furthermore, the sample size for measured flux densities limits a statistical inference on the general behaviour of minihalos. The final uncertainties for the flux densities measured for this thesis are dominated by the source-subtraction component, particularly for low surface-brightness minihalos.

A first step can be achieved by completing the self-calibration and source-subtraction for the calibrated datasets in this sample. This would enable a broader spectral comparison between the clusters. An exploration of the spatial dependence of the spectral index would also assist in distinguishing between the hadronic scenario and turbulent reacceleration. A spatial investigation comparing to lower frequency observations may reveal if these mechanisms operate over different minihalo regions, as previously observed for RXJ1720+2638. 

The findings of this thesis extend high-frequency constraints on minihalo spectra. The absence of widespread spectral steepening challenges simple reacceleration-only models. These observations constitute the first high-frequency survey of radio minihalos, highlighting the importance of multi-frequency and spatially resolved studies. 